\section*{Capítulo \ref{ch:b3:ocean-dynamics} --- Dinámica oceánica}

\begin{solution}{exo:b3:ocean-dynamics:1}
$f = \qty{9.37e-5}{s^{-1}}$; $\tau/(\rho f) = 0.12/(1025\times\num{9.37e-5})
= \qty{1.25}{m^2/s}$, a la derecha del viento: sur. A través de
\qty{1000}{km}, \qty{1.25e6}{m^3/s}, \qty{1.25}{\sverdrup}.
\end{solution}

\begin{solution}{exo:b3:ocean-dynamics:2}
$f = \qty{1.031e-4}{s^{-1}}$; $D_E = \pi\sqrt{2\times0.05/\num{1.031e-4}}
= \qty{98}{m}$. La corriente allí es $\mathrm e^{-\pi} = \num{0.043}$ de la
corriente superficial, que fluye en sentido contrario.
\end{solution}

\begin{solution}{exo:b3:ocean-dynamics:3}
$\beta = 2\Omega\cos\ang{30}/a = \qty{1.98e-11}{m^{-1}.s^{-1}}$; $V =
\num{-1e-7}/(1025\times\num{1.98e-11}) = \qty{-4.92}{m^2/s}$; a través de
\qty{5000}{km}, \qty{-24.6}{\sverdrup}: hacia el sur.
\end{solution}

\begin{solution}{exo:b3:ocean-dynamics:4}
Las dos \omterm{def:b3:ocean-dynamics:tides}{componentes de marea} vuelven a estar en fase tras $1/(1/12 - 1/12.4206) =
\qty{354}{h}$, \num{14.77}\text{ días}, la mitad de un mes sinódico.
\end{solution}

\begin{solution}{exo:b3:ocean-dynamics:5}
A \ang{40} N, $f = \qty{9.37e-5}{s^{-1}}$, $\beta =
\qty{1.75e-11}{m^{-1}.s^{-1}}$, $\tau_x = \qty{0.05}{Pa}$ y
$\partial\tau_x/\partial y = 0.1/\num{1.11e6} = \qty{9.0e-8}{Pa/m}$. El
primer término, $-\num{9.0e-8}/(1025\times\num{9.37e-5}) =
\qty{-9.38e-7}{m/s}$; el segundo, $0.05\times\num{1.75e-11}/(1025\times
(\num{9.37e-5})^2) = \qty{+9.7e-8}{m/s}$. $w_E = \qty{-8.4e-7}{m/s}$,
\qty{-26.5}{m/yr}: el agua es empujada hacia abajo.
\end{solution}

\begin{solution}{exo:b3:ocean-dynamics:6}
$\delta_S = \num{2e-6}/\num{2e-11} = \qty{100}{km}$; velocidad $=
\num{30e6}/(\num{1e5}\times800) = \qty{0.375}{m/s}$.
\end{solution}

\begin{solution}{exo:b3:ocean-dynamics:7}
$\beta = \qty{1.98e-11}{m^{-1}.s^{-1}}$; $\delta_M = (10^4/\num{1.98e-11})^{1/3}
= \qty{79.6}{km}$, y el chorro tiene $3.628\,\delta_M = \qty{289}{km}$ de anchura.
Con $A_H$ ocho veces menor, ambos se reducen a la mitad, ya que $8^{1/3} = 2$.
\end{solution}

\begin{solution}{exo:b3:ocean-dynamics:8}
$f = \qty{8.57e-5}{s^{-1}}$; $v = (21.3 - 14.5)/(\num{8.57e-5}\times\num{8e4})
= \qty{0.99}{m/s}$, positivo con $B$ a la derecha: mirando desde $A$ hacia $B$,
al este, la corriente se dirige hacia el norte. La superficie del mar se sitúa
$6.8/9.81 = \qty{0.69}{m}$ más alta en $B$; la columna de agua en $B$ es la
más ligera y más cálida, a la derecha de la corriente.
\end{solution}

\begin{solution}{exo:b3:ocean-dynamics:9}
Rango de marea viva $1.5K(1 + 0.46) = \qty{0.78}{m}$; rango de marea muerta $1.5K(1 - 0.46)
= \qty{0.29}{m}$; razón $1.46/0.54 = 2.7$.
\end{solution}

\begin{solution}{exo:b3:ocean-dynamics:10}
Igual que el \cref{thm:b3:ocean-dynamics:kelvin}, $u = \pm(g/c)\eta$ para $F(x \mp
ct)$, y $fu = -g\,\partial\eta/\partial y$ se convierte en $\beta_0y(\pm g/c)G =
-gG'$, así que $G = \exp(\mp\beta_0y^2/2c)$. Solo el signo superior, el de la
onda hacia el este, está acotada a ambos lados del ecuador. $\beta_0 =
\qty{2.29e-11}{m^{-1}.s^{-1}}$, $\sqrt{2.5/\num{2.29e-11}} = \qty{330}{km}$.
\end{solution}

\begin{solution}{exo:b3:ocean-dynamics:11}
(a)~Integrando $\beta v = f\,\partial w/\partial z$ desde el fondo, $w = 0$,
hasta la parte superior de la capa abisal, $w = W_0$: $\beta V = fW_0$, y $f/\beta =
a\tan\varphi$. (b)~$V = \num{6.371e6}\times0.577\times10^{-7} =
\qty{0.368}{m^2/s}$, \qty{1.84}{\sverdrup} a través de \qty{5000}{km}, hacia el polo.
(c)~Desde una corriente de frontera occidental profunda alimentada por las fuentes de alta latitud;
el interior fluye hacia las fuentes, y la corriente de frontera se aleja
de ellas.
\end{solution}

\begin{solution}{exo:b3:ocean-dynamics:12}
(a)~Con $z$ hacia arriba desde el fondo, $A_zW'' = \mathrm if(W - W_g)$, siendo
el gradiente de presión el del flujo geostrófico de encima. $W - W_g =
C\mathrm e^{-(1+\mathrm i)z/\delta}$ está acotado hacia arriba, y $W(0) = 0$ da
$C = -W_g$. (b)~El defecto de transporte es $\int_0^\infty(W - W_g)\,\mathrm dz
= -W_g\delta/(1+\mathrm i) = -\tfrac12\delta(1 - \mathrm i)W_g$, es decir,
$-\tfrac12\delta(u_g + v_g,\ v_g - u_g)$. Siendo el flujo geostrófico
no divergente, $w = -\nabla\cdot(\text{defecto}) = \tfrac12\delta(\partial_x
v_g - \partial_y u_g) = \tfrac12\delta\zeta_g$. (c)~$\delta =
\sqrt{0.02/10^{-4}} = \qty{14.1}{m}$; $r = 10^{-4}\times14.1/8000 =
\qty{1.77e-7}{s^{-1}}$, un tiempo de frenado de \num{65}\text{ días}.
\end{solution}

\begin{solution}{pb:b3:ocean-dynamics:1}
\textbf{1.} $\tau_x = \qty{-0.1}{Pa}$ (del este), 0 y \qty{+0.1}{Pa}
(del oeste).

\textbf{2.} $f = \num{4.99e-5}$, \num{8.37e-5} y
\qty{1.117e-4}{s^{-1}}.

\textbf{3.} $V_E = -\tau_x/(\rho f)$: \qty{+1.96}{m^2/s} (norte) a
\ang{20} N y \qty{-0.87}{m^2/s} (sur) a \ang{50} N. Convergen.

\textbf{4.} $D_E = \pi\sqrt{2\times0.05/\num{8.37e-5}} = \qty{109}{m}$.

\textbf{5.} $\tau_x = 0$ a \ang{35} N, por tanto $w_E = -(1/\rho f)\,
\partial\tau_x/\partial y = -\pi\tau_0/(\rho fL) = \qty{-1.10e-6}{m/s}$,
\qty{-34.7}{m/yr}, hacia abajo.

\textbf{6.} $500/34.7 = \num{14.4}\text{ años}$.

\textbf{7.} $L = a\times\pi/6 = \qty{3336}{km}$; $\beta =
\qty{1.875e-11}{m^{-1}.s^{-1}}$.

\textbf{8.} $\operatorname{curl}_z\vec\tau = -\pi\tau_0/L =
\qty{-9.42e-8}{Pa/m}$.

\textbf{9.} $V = \num{-9.42e-8}/(1025\times\num{1.875e-11}) =
\qty{-4.90}{m^2/s}$, hacia el sur.

\textbf{10.} $4.90\times\num{6e6} = \qty{2.94e7}{m^3/s}$,
\qty{29.4}{\sverdrup} hacia el sur.

\textbf{11.} $\qty{4.9}{mm/s}$; \qty{1000}{km} en \qty{2.0e8}{s},
\num{6.5}~años.

\textbf{12.} $fw_E/\beta = \num{8.37e-5}\times(\num{-1.10e-6})/
\num{1.875e-11} = \qty{-4.90}{m^2/s}$: lo mismo, porque el transporte
de Ekman se anula donde $\tau_x = 0$. Una columna que se aplasta a un ritmo $w_E$
conserva su \omterm{def:b3:atmospheric-dynamics:pv}{vorticidad potencial} $f/h$ solo desplazándose hacia un $f$ menor, hacia el sur,
a un ritmo $\beta V_g = fw_E$.

\textbf{13.} Una capa límite solo puede serlo si la corrección a la solución
del interior se atenúa al alejarse de la pared; con $\beta > 0$ eso es
$\mathrm e^{-x/\delta_S}$, que se atenúa hacia el este, desde la pared
occidental únicamente. En términos de vorticidad, el rozamiento en una corriente hacia el norte contra
una pared oriental eliminaría vorticidad positiva, con el signo equivocado.

\textbf{14.} $\delta_S = \num{e-6}/\num{1.875e-11} = \qty{53.3}{km}$.

\textbf{15.} $A = \qty{4.90}{m^2/s}$, $X/\delta_S = 112.5$:
$4.90\times111.5 = \qty{546}{m^2/s}$, \qty{0.55}{m/s} a lo largo de
\qty{1000}{m}.

\textbf{16.} Donde $(X/\delta_S)\mathrm e^{-x/\delta_S} = 1$: $x =
\delta_S\ln(X/\delta_S) = \qty{252}{km}$.

\textbf{17.} $A(X - \delta_S - x_0) = 4.90\times(6000 - 53 - 252)\times
10^3 = \qty{27.9}{\sverdrup}$. El flujo interior sigue dirigido hacia el sur en
la franja junto a la pared; el déficit, $A(x_0 + \delta_S) =
\qty{1.5}{\sverdrup}$, es el transporte interior que habría cruzado
esa franja.

\textbf{18.} $\delta_M = (\num{5e3}/\num{1.875e-11})^{1/3} = \qty{64.4}{km}$;
chorro $3.628\times64.4 = \qty{233}{km}$.

\textbf{19.} $1.163\times29.4 - A\times\qty{233}{km} = 34.2 - 1.1 =
\qty{33.0}{\sverdrup}$; el exceso, \qty{16}{\percent} del transporte
interior, vuelve hacia el sur en la contracorriente justo al este del chorro.

\textbf{20.} $c = \sqrt{9.81\times4000} = \qty{198}{m/s}$; $L_R =
198/\num{8.37e-5} = \qty{2370}{km}$.

\textbf{21.} $198\times12.4206\times3600 = \qty{8860}{km}$.

\textbf{22.} Perímetro $2(6000 + 3336) = \qty{18670}{km}$, \qty{26.2}{h},
2.1 periodos de M$_2$.

\textbf{23.} El círculo mide $2\pi a\cos\ang{35} = \qty{32790}{km}$, recorrido
en un día lunar, \qty{24.84}{h}: \qty{367}{m/s}, lo que requiere $c^2/g =
\qty{13.7}{km}$ de agua.

\textbf{24.} Con un radio de Rossby comparable a la anchura de la cuenca
y una longitud de onda de M$_2$ comparable a su longitud, aproximadamente un \omterm{def:b3:ocean-dynamics:amphidrome}{punto
anfidrómico}, cerca del centro de la cuenca; la pleamar gira en sentido antihorario
en torno a él, dejando la costa a su derecha.

\textbf{25.} $\tfrac34\times0.357\times\cos^2\ang{35} = \qty{0.18}{m}$.

\textbf{26.} La corriente de frontera occidental devuelve hacia el norte los
\textbf{\qty[text-series-to-math=true]{29}{\sverdrup}} que el interior
transporta hacia el sur (\qty{29.4}{\sverdrup} según el cálculo).
\end{solution}
